\section{Limitations and how to check them}
\label{sec:limitations}

The limitations below are grouped by where they come from, because the
grouping decides what to do about them: method-family limits will not
be removed by further development of this package, implementation
limits are measured and tracked, and protocol limits are the user's to
manage.  Each item comes with the check that detects it.

\subsection{Method-family limits}

\textbf{The perturbation layer inherits the standard contracted-NEVPT2
approximations.}  Excitation energies are reproduced at the accuracy
band the reference set assigns to \textsc{sc}-NEVPT2
(Chapter~\ref{sec:benchmark-walkthrough}), and not better.  Check:
compare against a second implementation or the reference set's own
NEVPT2 column before attributing a deviation to anything else.

\textbf{Valence--Rydberg mixing is not resolved at small bases.}
$\pi \to \pi^*$ states that sit near Rydberg states --- ethylene is the
textbook case --- shift by around an electron volt when diffuse
functions enter the basis, and move again when the active window
includes the diffuse pi orbitals.  Check: run the target once with and
once without diffuse functions; if the excitation moves by more than
the claimed accuracy, the number is basis-limited, not converged.

\textbf{Multiconfigurational $\pi\pi^*$ states exceed a two-orbital
window.}  The uniform positive deviation of the $\pi \to \pi^*$ family
in the benchmark chapter survives every window correction tried here.
Check: widen the window by one occupied orbital and re-run; a
deviation that shrinks is window-limited, one that does not is a
method-family limit for this protocol.

\textbf{No spin--orbit coupling and no relativistic Hamiltonian.}  The
package works in a scalar-relativistic world with ECPs; fine
structure and heavy-element spin physics are out of scope.  Check:
nothing to check --- if the target needs SOC, this is the wrong tool.

\subsection{Implementation limits}

\textbf{Dense diagonalization bounds the active space.}  The CI step
builds and diagonalizes the full determinant space; spaces beyond a
few thousand determinants become impractical, and iterative
eigensolvers are future work.

\textbf{Memory and compile time grow steeply with the parameter
count.}  The differentiated chain broadcasts intermediates over the
rotation parameters; at around five hundred parameters a single node
of the Hessian path once requested a 402~GB allocation, which is why
the second-order stages now run on Hessian-vector products.  Expect
the first evaluation of a large system to spend most of its time in
XLA compilation.  Check: watch the resident memory of the first
optimization step, not the steady state.

\textbf{CPU-only validation.}  The GPU path is not claimed; the
platform must be pinned to CPU (Chapter~\ref{sec:installation}).

\subsection{Protocol limits}

\textbf{The convergence gate is the gradient.}  Review
Chapter~\ref{sec:convergence} before quoting any number; the
stationary-point spread is part of the result.

\textbf{The window rules cover the patterns met so far.}  The frontier
audit of Chapter~\ref{sec:active-space} catches the two failure modes
this project hit; it is a rule set, not a guarantee, and a chemistry
outside the covered patterns should be audited with the same tools
rather than assumed correct.

\textbf{Residual deviations carry protocol, not just method.}  The
benchmark protocol fixes the basis at 6-31G and the geometry at the
reference set's; a number quoted without those two facts travels
badly.  Check: every table in this documentation states basis,
geometry source, window rule, and state-weighting pass --- keep that
habit when exporting numbers.
